\documentclass{article}
\usepackage[UTF8]{ctex}
\usepackage{amsmath,amssymb,amsfonts}
\usepackage{enumitem}
\usepackage{geometry}
\geometry{a4paper, margin=1in}
% === 错题本排版优化（不参与题目标识） ===
\setlength{\parindent}{0pt}
\setlength{\parskip}{0.35em plus 0.10em minus 0.08em}
\setlength{\emergencystretch}{3em}
\setlist{itemsep=0.25em, topsep=0.40em, parsep=0pt}
\allowdisplaybreaks[2]
% === 错题本排版优化结束 ===
\begin{document}
\section*{浙江大学2024年数学分析真题}
\section*{101．浙江大学2024年数学分析试卷}

\subsection{题目1}
1．计算题．

\begin{itemize}
  \item[（1）] 求极限
\[
\lim _{n \rightarrow \infty} \sqrt[n]{\cos \frac{\pi}{2 n} \cdot \cos \frac{2 \pi}{2 n} \cdot \cos \frac{(n-1) \pi}{2 n}}
\]
  \item[（2）] 求积分
\[
\iint_{x^{2}+y^{2} \leq 1}|3 x+4 y| \mathrm{d} x \mathrm{~d} y
\]
  \item[（3）] 求极限
\[
\lim _{x \rightarrow 0}\left[\frac{1}{\ln (1+x)}-\frac{1}{\ln \left(x+\sqrt{1+x^{2}}\right)}\right] .
\]
  \item[（4）] 求极限
\[
\lim _{x \rightarrow+\infty} \frac{1}{x} \int_{0}^{x}|\sin t| \mathrm{d} t
\]
\end{itemize}

\subsection{题目2}
2．解答如下问题：

\begin{itemize}
  \item[（1）] 请叙述确界存在定理．
  \item[（2）] 请利用闭区间套定理证明确界存在定理．
\end{itemize}

\subsection{题目3}
3．设 \(x_{1}=1, x_{n+1}=\sqrt{\frac{2 x_{n}^{2}}{2+x_{n}^{2}}}, n=1,2,3, \cdots\) ．

\begin{itemize}
  \item[（1）] 证明： \(\lim _{n \rightarrow \infty} x_{n}=0\) ．
  \item[（2）] 计算 \(\lim _{n \rightarrow \infty} n\left(x_{n}-\ln \left(1+x_{n}\right)\right)\) ．
\end{itemize}

\subsection{题目4}
4．设函数 \(f_{1}(x)\) 在 \([0,1]\) 上连续，\(K(x, y)\) 在 \([0,1] \times[0,1]\) 上连续．对任意 \(x \in[0,1]\) ，定义

\[
f_{n+1}(x)=\int_{0}^{x} K(x, y) f_{n}(y) \mathrm{d} y, n=1,2
\]

请证明：函数列 \(\left\{f_{n}(x)\right\}\) 在 \([0,1]\) 上一致收敛于 0 。

\subsection{题目5}
5．设函数 \(f, g\) 在 \([0,1]\) 上连续，且存在包含于 \([0,1]\) 的数列 \(\left\{x_{n}\right\}\) ，使得对于任意 \(n \geq 1\) ，有 \(f\left(x_{n}\right)=\) \(g\left(x_{n+1}\right)\) ．请证明：存在 \(\xi \in[0,1]\) ，使得 \(f(\xi)=g(\xi)\) ．

\subsection{题目6}
6．设 \(f\) 是定义在 \([a, b]\) 上的函数，且 \(G=\{(x, f(x)) \mid x \in[a, b]\}\) 是 \(\mathbb{R}^{2}\) 上的有界闭集．请证明：\(f\) 是\([a, b]\) 上的连续函数．

\subsection{题目7}
7．设函数 \(f\) 是 \((0,1]\) 上的无界的单调函数，且广义积分 \(\int_{0}^{1} f(x) \mathrm{d} x\) 收敛．证明：

\[
\lim _{n \rightarrow \infty} \frac{1}{n} \sum_{k=2}^{n} f\left(\frac{k-1}{n}\right)=\int_{0}^{1} f(x) \mathrm{d} x
\]

\subsection{题目8}
8．设函数 \(f\) 在 \(\mathbb{R}\) 上可微，且满足对任意 \(x \in \mathbb{R}\) ，有 \(f(x+1)-f(x)=f^{\prime}(x)\) 以及 \(\lim _{x \rightarrow+\infty} f^{\prime}(x)=1\) ．请证明：存在常数 \(C\) ，使得 \(f(x)=x+C\) 。

\subsection{题目9}
9．设函数 \(f\) 在 \([0,1]\) 上连续，定义

\[
g(t)=\int_{0}^{1} \frac{t f(x)}{x^{2}+t^{2}} \mathrm{~d} x, t \in \mathbb{R}
\]

请证明：函数 \(g\) 在点 0 处连续当且仅当 \(f(0)=0\) ．
\end{document}
